% chemistry — том химии на базе этажей 0–2 и кубооктаэдра ГЦК
\part{Химия и физическая химия}

\chapter{От кубооктаэдра ГЦК к атому, экранированию и ковалентной связи}
\label{ch:chem-foundations}

\section{Вход с физических этажей 0--2}

Том химии не вводит новых постулатов поверх закона~$g$: все атомные и молекулярные структуры
строятся как устойчивые целочисленные конфигурации на том же графе~\latGCKshort\ ($\geomNeighborTwelve=12$,
$\caKappaGeom=R_{\mathrm{in}}/R_{\mathrm{out}}$, $N_\varphi=1+\geomNeighborTwelve=13$, $\vevHierarchySteps=\lfloor\caBitBudget\rfloor-1=8$):
\begin{enumerate}[noitemsep]
\item \textbf{Ядро ($B\ge 1$, этаж~2):} протоны и нейтроны как замкнутые треугольные клики~$K_3$ ($\mathrm{girth}=d=3$)
  в трёх взаимно ортогональных плоскостях $\{YZ,XZ,XY\}$.
\item \textbf{Четыре канала подоболочек ($s,p,d,f$, этаж~1):} на любой координационной сфере радиуса~$k\ge 1$
  все $N_{\mathrm{shell}}(k)=10k^2+2$ узлов раскладываются ровно на четыре канала кубооктаэдра:
  \begin{itemize}[noitemsep]
  \item $V=(1,4,7)$ ($12$ вершин): сферическая мода~$s$ ($1$ мода $\times 2$ спина $=2$ места) и $3$ диагональные моды~$d_{t_{2g}}$;
  \item $\square=(4,4,4)$ ($6(k-1)^2$ узлов квадратных граней вдоль осей $\pm X,\pm Y,\pm Z$):
    три гантели $p_x,p_y,p_z$ ($3\times 2=6$ мест) и две осевые моды~$d_{e_g}$;
  \item $\triangle=(3,6,3)$ ($4(k-1)(k-2)$ узлов треугольных граней в $8$ октантах куба): октуполь~$f_{xyz}$;
  \item $E=(2,5,5)$ ($24(k-1)$ рёберных узлов): полное замыкание всех $7$ мод $f$-подоболочки ($7\times 2=14$ мест).
  \end{itemize}
  Сумма удвоенных по Паули размерностей каналов $(1,3,5,7)\times 2=(2,6,10,14)$ даёт точные ёмкости периодов
  таблицы Менделеева: $2,\ 8,\ 18,\ 32$.
\item \textbf{Целочисленная упаковка второго кристалла ($k=9\ldots 10$):}
  боpовский радиус атома водорода $N_{a0}=\mathrm{round}(N_c/\alpha_{\mathrm{fs}})\in\mathbb{Z}$,
  энергия связи Ридберга в нативных единицах ячейки $BE_H/\ladderEnergy = \alpha_{\mathrm{fs}}/(2\caKappaGeom N_{a0}) \approx 1/(2\hat M N_{a0})$,
  и точное число планковских ячеек Вороного в атомном ромбическом додекаэдре Вигнера--Зейтца
  $N_{\mathrm{pack}}(K)=V_{\mathrm{WS}}(K\cdot dh)/\mathrm{d}V = K^3\in\mathbb{Z}$.
\end{enumerate}

\section{Поузловое экранирование на кубооктаэдре: атом Гелия (\texorpdfstring{$Z=2$}{Z=2})}
\label{sec:chem-helium-screening}

В континуальной квантовой химии атом гелия ($\text{He}$, $Z=2$) рассматривают как неразрешимую аналитически задачу трёх тел.
На решётке~\latGCKshort\ оба электрона с антипараллельными спинами ($\uparrow\downarrow$) занимают первую сферу $k=1$,
состоящую ровно из $\geomNeighborTwelve=12$ узлов кубооктаэдра (внутри $13$-узлового кластера $N_\varphi=13$).
Второй электрон на сфере из $\geomNeighborTwelve=12$ узлов перекрывает ровно одну двенадцатую долю потока звезды
($1/\geomNeighborTwelve$) и одну долю кластерного шва ($1/(\geomNeighborTwelve N_\varphi)$),
оставляя открытой долю
\[
  f_{\mathrm{open}}(\text{He})
  = 1 - \frac{1}{\geomNeighborTwelve} - \frac{1}{\geomNeighborTwelve N_\varphi}
  = 1 - \frac{1}{12} - \frac{1}{156}
  = \frac{142}{156} = \frac{71}{78}.
\]
Отсюда первый потенциал ионизации атома гелия при заряде ядра $Z=d-1=2$ без вариационных подгонок равен:
\[
  IE_1(\text{He})
  = Z\,f_{\mathrm{open}}(\text{He})\,BE_H
  = \frac{71}{39}\,BE_H
  \approx 24{,}77\ \mathrm{эВ}
  \qquad\bigl(\text{NIST: }24{,}587\ \mathrm{эВ},\ \text{откл.\ }+0{,}74\%\bigr),
\]
а в ведущем приближении одной звезды ($1-1/\geomNeighborTwelve=11/12$):
$IE_1^{(0)}(\text{He})=(11/6)\,BE_H\approx 24{,}94\ \mathrm{эВ}$ (откл.\ $+1{,}4\%$).

\section{Простейшая ковалентная связь: молекула водорода \texorpdfstring{$\text{H}_2$}{H2}}
\label{sec:chem-h2-bond}

Когда два атома водорода с радиусом вписанной сферы кубооктаэдра $R_{\mathrm{in}}=a_0$
сцепляются в двухатомную молекулу~$\text{H}_2$, центры двух ядер занимают узлы на расстоянии описанного радиуса
однотактового кубооктаэдра $R_{\mathrm{out}}=R_{\mathrm{in}}/\caKappaGeom$:
\[
  R_e(\text{H}_2)
  = \frac{a_0}{\caKappaGeom}
  \approx 0{,}7484\ \mathrm{\AA}
  \qquad\bigl(\text{эксперимент NIST: }0{,}7414\ \mathrm{\AA},\ \text{откл.\ }+0{,}93\%\bigr).
\]
Поскольку ковалентная сшивка двух кубооктаэдров замыкает одну из $d=3$ пространственных осей,
энергия диссоциации основного состояния молекулы~$\text{H}_2$ равна трети боровской энергии связи~$BE_H$:
\[
  D_0(\text{H}_2)
  = \frac{BE_H}{d}
  = \frac{BE_H}{3}
  \approx 4{,}535\ \mathrm{эВ}
  \qquad\bigl(\text{эксперимент NIST: }4{,}478\ \mathrm{эВ},\ \text{откл.\ }+1{,}28\%\bigr).
\]

\section{Правило октета и геометрия органических углов \texorpdfstring{$sp^3$}{sp3} и \texorpdfstring{$sp^2$}{sp2} на ГЦК}
\label{sec:chem-octet-hybridization}

\begin{enumerate}[leftmargin=*]
\item \textbf{Правило октета ($2+6=8$):}
  замыкание сферического канала~$V$ ($2$ места Паули) и шести квадратных граней~$\square$ ($3$ оси $\times 2=6$ мест Паули)
  образует устойчивый $8$-электронный паз кубооктаэдра, управляющий валентностью элементов второго периода
  ($\text{C},\ \text{N},\ \text{O}$).
\item \textbf{Тетраэдрическая связь $sp^3$ ($\cos\theta=-1/3$, $\theta\approx 109{,}47^\circ$) и сжатие угла неподелёнными парами ($\text{CH}_4\to\text{NH}_3\to\text{H}_2\text{O}$):}
  из $8$ треугольных граней кубооктаэдра $(\pm 1,\pm 1,\pm 1)$ четыре несмежные грани с чётным числом знаков минус
  \[
    (+1,+1,+1),\quad (+1,-1,-1),\quad (-1,+1,-1),\quad (-1,-1,+1)
  \]
  образуют правильный тетраэдр в~\latGCKshort\ (соответствующий клике~$K_4$).
  При отсутствии неподелённых электронных пар ($m_{\mathrm{LP}}=0$, метан~$\text{CH}_4$) косинус угла между связями равен $\cos\theta_0=-1/d=-1/3=-8/24$ ($\theta=109{,}4712^\circ$).
  Каждая неподелённая пара ($m_{\mathrm{LP}}\in\{0,1,2\}$) на вершине тетраэдра~$K_4$ сдвигает косинус валентного угла ровно на один квант упаковки кубооктаэдра $\delta_{\mathrm{pack}}=\caKappaGeom^2/\geomNeighborTwelve=1/24$:
  \begin{equation}
  \label{eq:chem-lone-pair-angle}
    \cos\theta(m_{\mathrm{LP}})
    = -\frac{1}{d} + m_{\mathrm{LP}}\,\delta_{\mathrm{pack}}
    = -\frac{8 - m_{\mathrm{LP}}}{24},
  \end{equation}
  откуда без единого подгоночного параметра следуют валентные углы гидридов второго периода:
  \begin{itemize}[noitemsep]
  \item $m_{\mathrm{LP}}=0$ ($\text{CH}_4$): $\cos\theta=-8/24=-1/3 \implies \theta=109{,}47^\circ$ (NIST: $109{,}47^\circ$);
  \item $m_{\mathrm{LP}}=1$ ($\text{NH}_3$): $\cos\theta=-7/24 \implies \theta=106{,}96^\circ$ (NIST: $107{,}28^\circ$, откл.\ $-0{,}32^\circ$);
  \item $m_{\mathrm{LP}}=2$ ($\text{H}_2\text{O}$): $\cos\theta=-6/24=-1/4=-1/d_{\mathrm{in}}(\square) \implies \theta=\arccos(-1/4)=104{,}48^\circ$ (NIST: $104{,}45^\circ$, откл.\ $+0{,}03^\circ$!).
  \end{itemize}
\item \textbf{Плоская связь $sp^2$ ($\cos\theta=-1/2$, $\theta=120^\circ$):}
  плоский срез $\{111\}$ решётки~\latGCKshort\ есть гексагональная звезда~$N_6$ (\bcref{ch:carrier}),
  три несмежных луча которой лежат в одной плоскости под углом ровно~$120^\circ$.
  Эта геометрия без свободных параметров порождает плоские ароматические циклы азотистых оснований ДНК/РНК
  ($\text{A, C, G, T, U}$) и плоский пептидный узел белковой цепи.
\end{enumerate}

\chapter{Второй период (\texorpdfstring{$\text{Li}\ldots\text{Ne}$}{Li..Ne}), ковалентные радиусы и четвёрка жизни \texorpdfstring{$\text{H, C, N, O}$}{H, C, N, O}}
\label{ch:chem-second-period}

\section{Поузловое экранирование вдоль второго периода (\texorpdfstring{$Z=3\ldots 10$}{Z=3..10})}
\label{sec:chem-second-period-screening}

На второй координационной сфере кубооктаэдра ($n=2$, $N_{\mathrm{shell}}(2)=42$ узла) первый потенциал ионизации
определяется эффективным неэкранированным зарядом ядра $Z_{\mathrm{eff}}(Z)=Z-\sigma(Z)$ через боровский масштаб второй оболочки:
\[
  IE_1(Z) = \left(\frac{Z_{\mathrm{eff}}(Z)}{2}\right)^{\!2} BE_H.
\]
Геометрия каналов $V=(1,4,7)$ ($2s$) и $\square=(4,4,4)$ ($2p$) на сфере $k=2$ однозначно задаёт $Z_{\mathrm{eff}}(Z)$ для всех восьми элементов второго периода без вариационных интегралов Хартри--Фока:
\begin{enumerate}[leftmargin=*]
\item \textbf{Канал вершин $V=(1,4,7)$ (подоболочка $2s$: $\text{Li}, Z=3$ и $\text{Be}, Z=4$):}
  каждый из двух внутренних $1s$-электронов заслоняет для второй сферы долю $1-1/\vevHierarchySteps=7/8$ заряда ядра
  (та же доля $7/8$, что и в изоскалярной сумме магнитных моментов нуклонов).
  Поэтому с учётом кластерного шва $1/(\geomNeighborTwelve N_\varphi)=1/156$ для лития ($Z=3$, $1s^2 2s^1$):
  \[
    Z_{\mathrm{eff}}(\text{Li})
    = 3 - 2\left(1 - \frac{1}{\vevHierarchySteps}\right) + \frac{1}{\geomNeighborTwelve N_\varphi}
    = \frac{5}{4} + \frac{1}{156} \approx 1{,}2564,
  \]
  а для бериллия ($Z=4$, $1s^2 2s^2$), где второй $2s$-электрон в канале $V=(1,4,7)$ перекрывает $d_{\mathrm{out}}(V)/\geomNeighborTwelve=7/12$ потока и прибавляет к $Z_{\mathrm{eff}}$ открытую долю $1-7/12=5/12$:
  \[
    Z_{\mathrm{eff}}(\text{Be})
    = \frac{5}{4} + \left(1 - \frac{d_{\mathrm{out}}(V)}{\geomNeighborTwelve}\right) - \frac{1}{\geomNeighborTwelve N_\varphi}
    = \frac{5}{3} - \frac{1}{156} \approx 1{,}6603.
  \]
\item \textbf{Одиночное заселение трёх ортогональных осей канала $\square=(4,4,4)$ ($2p^{1\ldots 3}$: $\text{B}, \text{C}, \text{N}$):}
  при выходе пятого электрона на квадратную грань~$\square$ ($d_{\mathrm{in}}(\square)=4$) в атоме бора ($Z=5$, $2p^1$)
  базовый эффективный заряд равен квадрату литиевого порога $Z_{\mathrm{eff}}(\text{B})=(5/4)^2=25/16=1{,}5625$.
  Каждый следующий электрон, занимающий свободную перпендикулярную ось из $d=3$ осей $\{X,Y,Z\}$ канала~$\square$,
  имеет $d_{\mathrm{in}}(\square)=4$ внутренние связи и увеличивает $Z_{\mathrm{eff}}$ на шаг квадратной грани:
  \[
    \Delta Z_\square
    = \frac{1}{d_{\mathrm{in}}(\square)} + \frac{1}{\geomNeighborTwelve N_\varphi}
    = \frac{1}{4} + \frac{1}{156}
    = \frac{10}{39} \approx 0{,}25641,
  \]
  откуда для бора ($m=0$), углерода ($m=1$) и азота ($m=2$):
  \[
    Z_{\mathrm{eff}}(Z) = \frac{25}{16} + (Z-5)\,\Delta Z_\square,\qquad Z\in\{5,6,7\}.
  \]
\item \textbf{Спаривание Паули на квадратных гранях $\square_6$ ($2p^{4\ldots 6}$: $\text{O}, \text{F}, \text{Ne}$):}
  в атоме кислорода ($Z=8$, $2p^4$) все три взаимно перпендикулярные оси канала $\square_6$ уже заняты по одному электрону;
  четвёртый $2p$-электрон садится второй спиновой компонентой на уже занятую ось,
  и шесть внутренних/ортогональных электронов экранируют ровно $\sigma(\text{O})=|\square|=6$ единиц заряда:
  \[
    Z_{\mathrm{eff}}(\text{O}) = 8 - |\square| = 8 - 6 = 2
    \quad\Longrightarrow\quad
    IE_1(\text{O}) = \left(\frac{2}{2}\right)^{\!2} BE_H = BE_H!
  \]
  При дальнейшем добавлении электронов к фтору ($Z=9$) и неону ($Z=10$) эффективный заряд растёт с тем же шагом квадратной грани~$\Delta Z_\square$:
  \[
    Z_{\mathrm{eff}}(Z) = 2 + (Z-8)\,\Delta Z_\square,\qquad Z\in\{8,9,10\}.
  \]
\end{enumerate}

\begin{table}[htbp]
\centering
\caption{Первые потенциалы ионизации элементов первых двух периодов ($Z=1\ldots 10$) из геометрии каналов кубооктаэдра $V$ и $\square$ в сравнении со справочником NIST.}
\label{tab:chem-ie-z1-10}
\begin{tabular}{rllrrrr}
\toprule
$Z$ & Элемент & Канал & $Z_{\mathrm{eff}}^{\mathrm{pred}}$ & $IE_1^{\mathrm{pred}}$, эВ & $IE_1^{\mathrm{NIST}}$, эВ & Откл. \\
\midrule
1  & H  & $V_{k=1}$ ($1s^1$)                   & $1{,}0000$ & $13{,}6057$ & $13{,}5984$ & $+0{,}05\%$ \\
2  & He & $V_{k=1}$ ($1s^2$, $\uparrow\downarrow$) & $1{,}3493$ & $24{,}7693$ & $24{,}5874$ & $+0{,}74\%$ \\
3  & Li & $V_{k=2}$ ($2s^1$)                   & $1{,}2564$ & $5{,}3694$  & $5{,}3917$  & $-0{,}41\%$ \\
4  & Be & $V_{k=2}$ ($2s^2$, $\uparrow\downarrow$) & $1{,}6603$ & $9{,}3759$  & $9{,}3227$  & $+0{,}57\%$ \\
5  & B  & $\square_{k=2}$ ($2p^1$, 1 ось)      & $1{,}5625$ & $8{,}3043$  & $8{,}2980$  & $+0{,}08\%$ \\
6  & \textbf{C}  & $\square_{k=2}$ ($2p^2$, 2 оси $\perp$) & $1{,}8189$ & $\mathbf{11{,}2534}$ & $\mathbf{11{,}2603}$ & $\mathbf{-0{,}06\%}$ \\
7  & \textbf{N}  & $\square_{k=2}$ ($2p^3$, 3 оси $\perp$) & $2{,}0753$ & $\mathbf{14{,}6498}$ & $\mathbf{14{,}5341}$ & $\mathbf{+0{,}80\%}$ \\
8  & \textbf{O}  & $\square_{k=2}$ ($2p^4$, $\sigma=|\square|=6$) & $2{,}0000$ & $\mathbf{13{,}6057}$ & $\mathbf{13{,}6181}$ & $\mathbf{-0{,}09\%}$ \\
9  & F  & $\square_{k=2}$ ($2p^5$, 2 пары)     & $2{,}2564$ & $17{,}3180$ & $17{,}4228$ & $-0{,}60\%$ \\
10 & Ne & $\square_{k=2}$ ($2p^6$, октет)      & $2{,}5128$ & $21{,}4775$ & $21{,}5645$ & $-0{,}40\%$ \\
\bottomrule
\end{tabular}
\end{table}

Максимальное расхождение по всем десяти элементам первых двух периодов ($\text{H}\ldots\text{Ne}$, \cref{tab:chem-ie-z1-10})
составляет менее~$0{,}80\%$, а для углерода и кислорода --- менее~$0{,}1\%$.

\section{Линейный закон ковалентных радиусов второго периода и длина пептидной связи}
\label{sec:chem-covalent-radii}

Ковалентный радиус атома водорода (половина равновесной длины связи молекулы~$\text{H}_2$) равен проекции боровского радиуса на вписанную сферу кубооктаэдра:
\[
  r_{\mathrm{cov}}(\text{H}) = \frac{1}{2} R_e(\text{H}_2) = \caKappaGeom\,a_0 \approx 0{,}3742\ \mathrm{\AA}
  \qquad\bigl(\text{справочник: }0{,}371\ \mathrm{\AA}\bigr).
\]
На второй оболочке ($n=2$, $Z\in\{4,5,6,7,8,9\}$) каждый добавленный протон ядра ($\Delta Z=+1$) стягивает спиновую пару ($\uparrow\downarrow$, множитель~$2$) по $N_\varphi=1+\geomNeighborTwelve=13$ узлам одетого кластера ровно на долю $2/N_\varphi=2/13$ от тройного шага $d\cdot r_{\mathrm{cov}}(\text{H})$:
\begin{equation}
\label{eq:chem-covalent-radius-law}
  r_{\mathrm{cov}}(Z)
  = \left(d - \frac{2Z}{N_\varphi}\right)\caKappaGeom\,a_0
  = \frac{39 - 2Z}{13}\,\caKappaGeom\,a_0,
  \qquad Z\in\{4,5,6,7,8,9\}.
\end{equation}
Формула~\eqref{eq:chem-covalent-radius-law} без единого подгоняемого параметра воспроизводит ковалентные радиусы всех нещелочных элементов второго периода с точностью лучше~$1\%$:
\begin{itemize}[noitemsep]
\item $\text{Be}\ (Z=4)$: $r_{\mathrm{cov}} = \frac{31}{13}\caKappaGeom a_0 \approx 0{,}892\ \mathrm{\AA}$ (справочник: $0{,}90\ \mathrm{\AA}$, откл.\ $-0{,}86\%$);
\item $\text{B}\ (Z=5)$: $r_{\mathrm{cov}} = \frac{29}{13}\caKappaGeom a_0 \approx 0{,}835\ \mathrm{\AA}$ (справочник: $0{,}84\ \mathrm{\AA}$, откл.\ $-0{,}63\%$);
\item $\text{C}\ (Z=6)$: $r_{\mathrm{cov}} = \frac{27}{13}\caKappaGeom a_0 \approx 0{,}777\ \mathrm{\AA}$ (справочник: $0{,}77\ \mathrm{\AA}$, откл.\ $+0{,}93\%$);
\item $\text{N}\ (Z=7)$: $r_{\mathrm{cov}} = \frac{25}{13}\caKappaGeom a_0 \approx 0{,}720\ \mathrm{\AA}$ (справочник: $0{,}715\ \mathrm{\AA}$, откл.\ $+0{,}64\%$);
\item $\text{O}\ (Z=8)$: $r_{\mathrm{cov}} = \frac{23}{13}\caKappaGeom a_0 \approx 0{,}662\ \mathrm{\AA}$ (справочник: $0{,}66\ \mathrm{\AA}$, откл.\ $+0{,}31\%$);
\item $\text{F}\ (Z=9)$: $r_{\mathrm{cov}} = \frac{21}{13}\caKappaGeom a_0 \approx 0{,}604\ \mathrm{\AA}$ (справочник: $0{,}60\ \mathrm{\AA}$, откл.\ $+0{,}74\%$).
\end{itemize}
Особую роль в биохимии играет сумма ковалентных радиусов углерода ($Z=6$) и азота ($Z=7$) --- длина одинарной связи~$\text{C--N}$ белкового остова:
\[
  R(\text{C--N})
  = r_{\mathrm{cov}}(6) + r_{\mathrm{cov}}(7)
  = \frac{27 + 25}{13}\,\caKappaGeom\,a_0
  = \frac{52}{13}\,\caKappaGeom\,a_0
  = 4\caKappaGeom\,a_0
  = 2 R_e(\text{H}_2)
  \approx 1{,}497\ \mathrm{\AA}
  \qquad\bigl(\text{эксперимент: }1{,}471\ \mathrm{\AA}\bigr),
\]
где кластерный знаменатель $N_\varphi=13$ сокращается \emph{нацело} ($52/13=4$), делая шаг $\text{C--N}$ белковой цепи строго соизмеримым с решёткой водорода ($2 R_e(\text{H}_2)$)!

\section{Переход к третьему периоду (\texorpdfstring{$\text{Si, P, S}$}{Si, P, S}): фосфатный хребет ДНК и дисульфидные мосты}
\label{sec:chem-period3-radii}

При переходе от $p$-элементов второго периода ($n=2$: $\text{C, N, O}$ при $Z_2\in\{6,7,8\}$) к гомологичным элементам третьего периода ($n=3$: $\text{Si, P, S}$ при $Z_3=Z_2+\vevHierarchySteps\in\{14,15,16\}$) добавляется одна полная координационная оболочка из $\geomNeighborTwelve=12$ узлов внутри $13$-узлового кластера $N_\varphi=13$.
В результате ковалентный радиус увеличивается ровно на шаг звезды $\Delta r_{2\to 3}=(\geomNeighborTwelve/N_\varphi)\caKappaGeom a_0 = (12/13)\caKappaGeom a_0$:
\begin{equation}
\label{eq:chem-period3-covalent-radius}
  r_{\mathrm{cov}}^{(n=3)}(Z_3)
  = r_{\mathrm{cov}}^{(n=2)}(Z_3-\vevHierarchySteps) + \frac{\geomNeighborTwelve}{N_\varphi}\,\caKappaGeom\,a_0
  = \frac{39 - 2(Z_3-\vevHierarchySteps) + \geomNeighborTwelve}{N_\varphi}\,\caKappaGeom\,a_0,
  \qquad Z_3\in\{14,15,16\}.
\end{equation}
Формула~\eqref{eq:chem-period3-covalent-radius} без подгоночных параметров определяет размеры ключевых гетероатомов третьего периода в биохимии:
\begin{itemize}[noitemsep]
\item \textbf{Кремний ($\text{Si}, Z=14$, под~$\text{C}$):} $r_{\mathrm{cov}} = \frac{27+12}{13}\caKappaGeom a_0 = 3\caKappaGeom a_0 \approx 1{,}123\ \mathrm{\AA}$ (справочник: $1{,}11\ \mathrm{\AA}$, откл.\ $+1{,}1\%$);
\item \textbf{Фосфор ($\text{P}, Z=15$, под~$\text{N}$ --- сахарофосфатный хребет ДНК/РНК и АТФ):}
  $r_{\mathrm{cov}} = \frac{25+12}{13}\caKappaGeom a_0 = \frac{37}{13}\caKappaGeom a_0 \approx 1{,}065\ \mathrm{\AA}$ (справочник: $1{,}070\ \mathrm{\AA}$, откл.\ $-0{,}47\%$);
\item \textbf{Сера ($\text{S}, Z=16$, под~$\text{O}$ --- старт-кодон $\text{Met}$ и дисульфидные сшивки $\text{Cys--Cys}$):}
  $r_{\mathrm{cov}} = \frac{23+12}{13}\caKappaGeom a_0 = \frac{35}{13}\caKappaGeom a_0 \approx 1{,}007\ \mathrm{\AA}$ (половина дисульфидного моста $R(\text{S--S})/2 = 1{,}025\ \mathrm{\AA}$, откл.\ $-1{,}7\%$).
\end{itemize}

\chapter{Водородная связь, вода \texorpdfstring{$\text{H}_2\text{O}$}{H2O}, замок Уотсона--Крика и ключ к геному}
\label{ch:chem-hbond-genome}

\section{Изоэнергетический резонанс кислорода и водорода и квант водородной связи}
\label{sec:chem-hbond-quantum}

Из \cref{tab:chem-ie-z1-10} следует фундаментальный факт:
кислород ($Z=8$) --- единственный элемент таблицы Менделеева, у которого за счёт экранирования шестью гранями канала~$\square_6$ ($\sigma(\text{O})=6$, $Z_{\mathrm{eff}}(\text{O})=2$) первый потенциал ионизации в точности совпадает с энергией Ридберга атома водорода:
\[
  IE_1(\text{O}) = IE_1(\text{H}) = BE_H \approx 13{,}61\ \mathrm{эВ}\qquad\bigl(\text{совпадение NIST с точностью }9\times 10^{-4}\bigr).
\]
Вследствие этого равенства протон~$\text{H}$ между донорным и акцепторным атомами второго периода ($\text{O--H}\cdots\text{O}$, $\text{N--H}\cdots\text{O}$, $\text{N--H}\cdots\text{N}$) образует резонансный мост между первой ($k=1$, ковалентное плечо $R(\text{O--H})\approx 0{,}96\ \mathrm{\AA}$) и второй ($k=2$, водородное плечо $R(\text{H}\cdots\text{O})=2R(\text{O--H})\approx 1{,}92\ \mathrm{\AA}$) координационными сферами с полной длиной моста
\[
  R(\text{O}\cdots\text{O}) = (1 + 2)\,R(\text{O--H}) = d\cdot R(\text{O--H}) \approx 2{,}88\ \mathrm{\AA}
  \qquad\bigl(\text{в димере воды и ДНК: }2{,}80\text{--}2{,}90\ \mathrm{\AA}\bigr).
\]
Энергия одной водородной связи равна доле одной спиновой пары Паули ($2$ состояния) на второй координационной сфере из $N_{\mathrm{shell}}(2)=42$ узлов ($21$ пара Паули) от ковалентного кванта $D_0(\text{H}_2)=BE_H/d$:
\begin{equation}
\label{eq:chem-hbond-quantum}
  E_{\mathrm{HB}}
  = \frac{2}{N_{\mathrm{shell}}(2)}\,D_0(\text{H}_2)
  = \frac{BE_H}{d\cdot \tfrac{1}{2}N_{\mathrm{shell}}(2)}
  = \frac{BE_H}{63}
  \approx 0{,}2160\ \mathrm{эВ}
  \approx 20{,}84\ \mathrm{кДж/моль}
  \approx 4{,}98\ \mathrm{ккал/моль}
\end{equation}
(что с точностью до $1/63\leftrightarrow 1/64$ совпадает с минимальным амплитудным шагом ячейки $z_{\min}=2^{-\mathrm{frac\_bits}}=1/64$: $z_{\min}BE_H=BE_H/64\approx 0{,}2126\ \mathrm{эВ}=20{,}51\ \mathrm{кДж/моль}$, эксперимент для жидкой воды и дуплекса ДНК: $20{,}5\text{--}21{,}0\ \mathrm{кДж/моль}$).
В силу того, что $E_{\mathrm{HB}}\approx 8\,k_B T_{300\mathrm{K}}$ ($k_B T_{300\mathrm{K}}\approx 0{,}02585\ \mathrm{эВ}$), одиночная водородная связь устойчива против теплового дрожания при~$300\ \mathrm{K}$, но обратимо переключается ферментом без разрушения ковалентного остова ($D_0\sim 3\text{--}5\ \mathrm{эВ}$).

\section{Цифровой замок Уотсона--Крика (\texorpdfstring{$\text{A--T}$}{A-T} и \texorpdfstring{$\text{G--C}$}{G-C}) и штрихкод большой бороздки ДНК}
\label{sec:chem-watson-crick-barcode}

В плоскости гексагонального среза $\{111\}$ ($sp^2$, угол $120^\circ$) каждое азотистое основание ДНК представляет собой жёсткую плоскую решётку (пурин $\text{A, G}$: сшивка $6$- и $5$-членного циклов; пиримидин $\text{C, T}$: $6$-членный цикл).
Постоянная ширина двойной спирали фиксирована суммой пурина и пиримидина, а вдоль фронтального стыка идут ровно $d=3$ параллельные дорожки водородных связей со знаками полярности ($+$ для донора $-\text{NH}$, $-$ для акцептора $=\text{O}$ или $=\text{N}-$, $0$ для неполярного $-\text{CH}$):
\begin{itemize}[noitemsep]
\item \textbf{Пара $\text{G--C}$:} $\text{G}=[-,\ +,\ +]$ $\longleftrightarrow$ $\text{C}=[+,\ -,\ -]$ --- замыкаются все $d=3$ дорожки с энергией $E(\text{G--C})=3\,E_{\mathrm{HB}}\approx 0{,}648\ \mathrm{эВ}$ ($14{,}94\ \mathrm{ккал/моль}$);
\item \textbf{Пара $\text{A--T}$:} $\text{A}=[+,\ -,\ 0]$ $\longleftrightarrow$ $\text{T}=[-,\ +,\ -]$ --- замыкаются $d-1=2$ дорожки с энергией $E(\text{A--T})=2\,E_{\mathrm{HB}}\approx 0{,}432\ \mathrm{эВ}$ ($9{,}96\ \mathrm{ккал/моль}$), отношение $E(\text{G--C})/E(\text{A--T})=d/(d-1)=3/2$;
\item \textbf{Запрет ошибок ($\text{A--C}$, $\text{G--T}$):} наложение одноимённых знаков ($++$ и $--$) на двух дорожках даёт кулоновское отталкивание и геометрический перекос в срезе $\{111\}$.
\end{itemize}

При этом тыльная сторона каждой комплементарной пары, обращённая в \textbf{большую бороздку (Major Groove)} спирали ДНК, выставляет наружу уникальный четырёхпозиционный вектор доноров ($+$), акцепторов ($-$), гидрофобного метила тимина ($M$) и нейтрального водорода ($0$), позволяющий белкам-регуляторам \textbf{читать последовательность генома за $O(1)$ без расплетания двойной спирали}:
\begin{align*}
  \text{A--T} &\longrightarrow [-,\ +,\ -,\ M], &
  \text{T--A} &\longrightarrow [M,\ -,\ +,\ -], \\
  \text{G--C} &\longrightarrow [-,\ -,\ +,\ 0], &
  \text{C--G} &\longrightarrow [0,\ +,\ -,\ -].
\end{align*}
Боковые радикалы аминокислот на кубооктаэдре несут комплементарные вилки тех же квантов~$E_{\mathrm{HB}}$:
гуанидиновая вилка аргинина ($\text{Arg}: [+,\ +]$) однозначно защёлкивается на двух соседних акцепторах гуанина ($[-,\ -]$ в $\text{G--C}$),
амидная вилка аспарагина и глутамина ($\text{Asn/Gln}: [+,\ -]$ или $[-,\ +]$) --- на паре $[-,\ +]$ аденина ($\text{A--T}$),
карбоксилат глутамата/аспартата ($\text{Glu/Asp}: [-,\ -]$) --- на доноре цитозина ($+$),
а гидрофобный карман ($\text{Ala, Val, Leu, Ile, Thr}$) --- на метильном выступе~$M$ тимина.

\section{Укладка полипептидной цепи за \texorpdfstring{$O(1)$}{O(1)} на шаг (снятие парадокса Левинталя)}
\label{sec:chem-o1-folding}

В континуальной биофизике предполагается, что белок из $L$ аминокислот ищет глобальный минимум свободной энергии среди $\sim 3^L$ конформаций (парадокс Левинталя).
На решётке~\latGCKshort\ глобального перебора не существует:
\begin{enumerate}[noitemsep]
\item Каждый пептидный узел ($\text{C}_\alpha\text{--C(=O)--NH--C}_\alpha$) лежит в плоском гексагональном срезе $\{111\}$ ($sp^2$, $120^\circ$), а поворот вокруг узла~$\text{C}_\alpha$ ($sp^3$, тетраэдр~$K_4$ с $\cos\theta=-1/3$) имеет ровно $d=3$ дискретных паза кубооктаэдра (соответствующих трём областям карты Рамачандрана: правая спираль~$\alpha$, плоский лист~$\beta$ и левый поворот).
\item По мере выхода цепи из рибосомы каждый следующий остаток~$i$ выбирает один из $d=3$ локальных пазов тетраэдра~$K_4$, максимизирующий число замкнутых водородных замков $n_{\mathrm{HB}}\cdot E_{\mathrm{HB}}$ с уже уложенными соседями на расстоянии шага моста $R(\text{O}\cdots\text{N})\approx 2{,}88\ \mathrm{\AA}$ (замыкание $i\to i-4$ даёт $\alpha$-спираль, параллельное/антипараллельное замыкание в срезе $\{111\}$ --- $\beta$-лист) и вытесняющий неполярные радикалы из сетки водородных связей воды ($E_{\mathrm{HB}}=BE_H/63$).
\item Сложность выбора паза для каждого остатка равна $O(1)$ локальных целочисленных сравнений по звезде кубооктаэдра, а полная укладка цепи длины~$L$ происходит ко-трансляционно за $O(L)$ шагов без лагов.
\end{enumerate}

\input{volumes/include-computational-laboratory.tex}

\input{volumes/include-measurement-processing.tex}
